\section{Query-Triggered Anomaly Detection}


To ensure the system remains responsive and memory-efficient, the transformation and detection logic are not pre-calculated. Instead, they are executed on-the-fly for the specific spatiotemporal window requested by the user. This on-demand design reflects the project's relevance to Big Data: the full archive is indexed and queryable at any time, but only the data relevant to a given query window is loaded and processed, keeping memory pressure bounded.

\subsection{Historical Baseline Construction}

All data with acquisition dates strictly before the user-configured evaluation start date is designated as the baseline. For each unique stable pixel ID and calendar month, the pipeline computes the baseline mean (\texttt{ndvi\_mu}, \texttt{ndmi\_mu}) and standard deviation (\texttt{ndvi\_sigma}, \texttt{ndmi\_sigma}) using Spark population standard deviation. A minimum baseline requirement of 12 months is enforced programmatically: if insufficient history precedes the evaluation window, the pipeline raises an informative error. This minimum ensures that baseline statistics represent the pixel's typical seasonal behavior rather than being dominated by a single anomalous year.

\subsection{Z-Score Computation and Anomaly Flagging}
\label{subsec:zscore}

For each pixel-month in the evaluation window, the observed NDVI (and NDMI) is compared to its historical monthly baseline using a standardized Z-score:

\begin{equation}
Z = \frac{\text{Observed Value} - \mu_{\text{monthly}}}{\sigma_{\text{monthly}}}
\end{equation}

where $\mu_{\text{monthly}}$ and $\sigma_{\text{monthly}}$ are the historical mean and standard deviation for that pixel and calendar month, as computed in the baseline phase. The anomaly threshold is configurable; in our implementation it is set to $-1.5$ rather than the more conservative $-2.0$ threshold, prioritizing recall (sensitivity to true stress events) over precision. A pixel-month with $Z < -1.5$ is flagged with \texttt{is\_ndvi\_anomaly = True} or \texttt{is\_ndmi\_anomaly = True}, respectively. Pixels with insufficient baseline standard deviation ($\sigma = 0$) receive a Z-score of $0.0$ and are not flagged. 

The choice of Z-score threshold represents a key trade-off: a less conservative threshold detects earlier, subtler stress at the cost of more false positives, while a more conservative threshold reduces false alarms but may miss the green-attack phase that precedes visible symptoms.

\subsection{Output Products}

Detected anomalies are exported to S3 in Apache Parquet format under a \texttt{run\_id} timestamp prefix (formatted as \texttt{YYYYMMDD\_HHMMSS}) for versioning and reproducibility. Two primary output datasets are produced per pipeline run: (1) an \texttt{anomaly\_events} Parquet file containing all flagged pixel-months with their Z-scores, vegetation index values, and baseline statistics; and (2) a \texttt{time\_series\_stats} Parquet file containing the complete pixel-month time series for post-hoc analysis and visualization. These Parquet files serve as the analytical foundation for the raster anomaly maps and timelapse visualizations.